% =====================================================================
% Chapter 5 — Route B: Single-View Learned Completion
% =====================================================================
\chapter{Route~B: Single-View Learned Completion}
\label{ch:routeB}

This chapter presents the methodology of \routeB, the learned route. As in the
previous chapter it first justifies \emph{why} Route~B is the appropriate
answer in this project (\secref{sec:why_routeB}), then details the pipeline in
its stages: single-view capture and sonar-realistic simulation
(\secref{sec:routeB_sim}), the PCN architecture (\secref{sec:routeB_arch}),
the construction of paired training data (\secref{sec:routeB_data}), and the
training procedure (\secref{sec:routeB_training}).

\section{Why Route~B?}
\label{sec:why_routeB}

Route~B is the answer to the research question under the following conditions:

\begin{enumerate}
  \item \textbf{Only one partial view is available (or will be at
  deployment).} The \suop{} acquisition protocol moved the object between
  cases without logging its pose, so the multi-view assumption of \routeA{}
  is not always satisfiable. In many operational settings --- a single sonar
  pass over an object, a towed sonar that sees an object once --- only one
  partial observation will ever exist. No geometric fusion can recover the
  occluded side of a single scan, so the only route to a \emph{complete}
  representation is a learned prior.
  \item \textbf{The output must be a complete object representation.}
  Route~A deliberately leaves unobserved regions empty; \routeB{} deliberately
  fills them. Applications that need a full shape --- e.g.\ a matching target
  for a manipulation system, a CAD-like model for simulation, or a
  classification feature --- prefer a complete, prior-informed output over a
  partial but strictly measured one.
  \item \textbf{A shape prior can be obtained.} Because the objects are
  manufactured or quasi-rigid (folding chairs, tyres, drums, mannequins,
  nets), their complete geometry is learnable: procedural parametric models and
  the project's own \routeA{} meshes provide paired training targets
  (\secref{sec:routeB_data}).
\end{enumerate}

The \emph{core assumption} of Route~B is that a network trained on many
(partial, complete) pairs learns a prior $p(\text{complete}\mid\text{partial})$
over the shape category, so that at inference a single partial cloud can be
mapped to a complete point set. Route~B is therefore a conditional-generation
problem: encode the partial view, decode a complete cloud, and supervise with a
symmetric distance to a reference. The route's pipeline is summarised in
\figref{fig:routeB_overview}.

\begin{figure}[H]
  \centering
  \includegraphics[width=0.9\textwidth]{fig_routeB_dummy.png}
  \caption{Route~B output for the dummy after 10 epochs: left, the sonar-like
  partial input; centre, the complete reference used as ground truth; right,
  the PCN fine completion. The completion contains inferred points in regions
  that were never visible in the input.}
  \label{fig:routeB_overview}
\end{figure}

\section{Stage 1: single-view capture and simulation}
\label{sec:routeB_sim}

\subsection{Real capture}
At inference, the input to \routeB{} is a single segmented sonar partial
cloud. In this project the real \suop{} chair segment shown in
\figref{fig:chair_pcn} is used directly: a single case's segmented returns are
normalised to the PCN input size of 2048 points and fed to the trained network
without any multi-view processing.

\begin{figure}[H]
  \centering
  \includegraphics[width=0.55\textwidth]{fig_chair_ping.png}
  \caption{A single real \suop{} chair ping (fan-shaped range return). A
  single view of this kind is the deployment input for \routeB{}: the object's
  far side, underside and much of its frame are simply absent from the data.}
  \label{fig:chair_pcn}
\end{figure}

\subsection{Simulated acquisition}
Because \suop{} contains no (partial, complete) ground-truth pairs, training
partials are generated from complete references. For every training sample the
following pipeline is applied (\texttt{simulate\_sonar\_partial} in
\texttt{train\_pcn.py}):

\begin{enumerate}
  \item \textbf{Viewpoint sampling.} A virtual camera is placed at a random
  azimuth $\phi \in [0, 2\pi)$ and elevation $\psi \in [-0.3, 0.6]$~rad at a
  distance of $2.5$--$4.0$ $\times$ the object diameter, i.e.\ representative of
  a 3~m-range sonar look at the published object scales.
  \item \textbf{Hidden-point removal.} Open3D's
  \texttt{hidden\_point\_removal} \citep{katz2007direct} keeps exactly the
  subset of the complete surface visible from that viewpoint, mimicking
  acoustic shadowing: the far side of the object is removed.
  \item \textbf{Sonar degradation.} Gaussian range noise with standard
  deviation $6$--$15$~mm is added to simulate beam/range error, and with
  probability $0.5$ a contiguous angular sector (span $0.5$--$1.5$~rad) is
  removed, imitating an occlusion shadow behind a strong return.
  \item \textbf{Resampling.} The resulting partial is (sub- or over-)sampled to
  exactly 2048 points, the PCN input size.
\end{enumerate}

The partial is normalised using the \emph{full shape} frame (centre and max
norm of the complete reference), so that the partial and the ground truth share
a common scale --- this is essential, because the Chamfer loss is
scale-sensitive. \figref{fig:sim_partial} illustrates the simulation.

\begin{figure}[H]
  \centering
  \includegraphics[width=0.9\textwidth]{fig_sim_partial.png}
  \caption{Sonar-realistic partial simulation. A complete reference shape is
  viewed from a random viewpoint (hidden-point removal keeps only visible
  points), then corrupted with Gaussian range noise and an angular sector
  removal that imitates an acoustic shadow.}
  \label{fig:sim_partial}
\end{figure}

\section{Stage 2: the PCN architecture}
\label{sec:routeB_arch}

The network is the Point Completion Network \citep{yuan2018pcn}, implemented
from scratch in PyTorch \citep{paszke2019pytorch}. Its design is summarised in
\tabref{tab:pcn_arch}.

\begin{table}[H]
  \centering
  \small
  \begin{adjustbox}{max width=\textwidth}
  \begin{tabular}{@{}llll@{}}
    \toprule
    Block & Operation & Output size & Notes \\
    \midrule
    Input & --- & $N \times 3$ & $N = 2048$ partial points \\
    Encoder MLP 1 & Conv1d($3\to128\to256$), ReLU & $256 \times N$ & shared per point \\
    Max pool (local) & \texttt{max} over points & $256 \times 1$ & PointNet symmetry \\
    Encoder MLP 2 & Conv1d($512\to512\to1024$), ReLU & $1024 \times N$ & concat with pooled \\
    Max pool (global) & \texttt{max} over points & $1024$ & global latent $\mathbf{v}$ \\
    FC decoder & Linear($1024\!\to\!1024\!\to\!1024\!\to\!3M$) & $M \times 3$ & coarse, $M = 1024$ \\
    Folding decoder & Conv1d($1024{+}2{+}3 \to 512 \to 3$) & $M g^2 \times 3$ & fine, $g \in \{2,4\}$ \\
    \bottomrule
  \end{tabular}
  \end{adjustbox}
  \caption{\pcn{} architecture used in this project. The encoder is a
  PointNet-style shared-MLP plus symmetric max pooling that produces a
  permutation-invariant 1024-dimensional global latent; the decoder predicts
  $M$ coarse points with a fully connected head and then folds a $g\times g$
  grid around each coarse point to produce $M g^2$ fine points.}
  \label{tab:pcn_arch}
\end{table}

\subsection{Encoder: why shared MLPs and max pooling}
Points are unordered, so any encoder must be \emph{permutation invariant}.
Shared MLPs applied per point (implemented as $1\times1$ convolutions) give a
feature for every point; the symmetric max pooling operation then aggregates
the strongest learned response per channel into a global feature that does not
depend on the order of the input points. The second MLP block sees the
concatenation of per-point features and the pooled feature, which gives each
point local context before the final pooling step. The output is a single
1024-dimensional vector summarising the partial shape.

\subsection{Decoder: why coarse-to-fine with folding}
Decoding 4096 structured points from one vector is a difficult regression. The
coarse-to-fine strategy splits it: a fully connected head predicts $M=1024$
coarse points that capture the global skeleton (seat, legs, rim, barrel, etc.);
each coarse point then \emph{sprouts} a $g\times g$ grid of fine points via
the folding operation
\begin{equation}
  \mathbf{f}_{i,k} \;=\; \mathbf{c}_i \;+\;
  \mathrm{MLP}\big( [\,\mathbf{c}_i;\; \mathbf{u}_{i,k};\; \mathbf{v}\,]\big),
  \qquad \mathbf{u}_{i,k} \in [-0.05,\;0.05]^2 ,
  \label{eq:folding}
\end{equation}
where a per-point MLP maps the
concatenation of the coarse point, the grid coordinate and the global latent to
a small offset. Folding lets the network refine local detail while remaining
tied to the global latent, which is why it is preferred over a second fully
connected head \citep{yang2018folding,yuan2018pcn}. With $g=2$ the fine cloud
has 4096 points; with $g=4$, 16{,}384 points. The output is a point cloud,
never a mesh.

\section{Stage 3: paired training data}
\label{sec:routeB_data}

\subsection{Two generators}
\label{sec:routeB_generators}
Complete reference clouds come from two generators:

\begin{enumerate}
  \item \textbf{Generator 1 --- procedural parametric models.}
  \texttt{make\_chair}, \texttt{make\_tire}, \texttt{make\_drum} and
  \texttt{make\_dummy} in \texttt{train\_pcn.py} build Open3D meshes from
  randomised dimensions matched to the published \suop{} nominal sizes: a
  folding chair from boxes (seat $0.40$--$0.50$~m, backrest, four legs
  $0.03$--$0.05$~m thick), a tyre from a torus ($R = 0.28$--$0.34$~m,
  $r = 0.08$--$0.13$~m), a drum from a cylinder (radius $0.26$--$0.34$~m,
  height $1.0$--$1.3$~m), and a dummy mannequin from a torso cylinder, head
  sphere, and limb cylinders with randomised proportions. Each mesh is sampled
  uniformly to 8192 surface points.
  \item \textbf{Generator 2 --- mesh-conditioned on \routeA{} output.}
  \texttt{gen\_suop\_dummy\_dataset.py} builds training pairs from the
  \emph{actual reconstructed \routeA{} meshes} (e.g.\ \texttt{meshes/dummy\_bpa.ply}):
  it samples 100{,}000 surface points from the mesh, then for each sample
  applies a random yaw rotation and a $\pm4\%$ scale jitter before running the
  partial simulation of \secref{sec:routeB_sim}. This deliberately links the
  two routes: the learned completion prior is conditioned on the geometry that
  the geometric route actually recovered from sonar, rather than on an ideal
  CAD shape.
\end{enumerate}

\subsection{Output pairs}
Each training sample is a pair $(\mathbf{P}, \mathbf{G})$: a 2048-point
sonar-like partial cloud and a 4096-point complete reference, both normalised
to the unit sphere in the full-shape frame. Datasets were generated with
$n = 6000$ pairs (e.g.\ \texttt{dataset\_dummy\_suop.npz},
\texttt{dataset\_drum\_suop.npz}) using a multiprocessing pool, in a few
minutes on eight cores. A 10\% hold-out provides validation.

\section{Stage 4: training}
\label{sec:routeB_training}

\subsection{Loss}
The training objective is the symmetric Chamfer distance
\eqref{eq:chamfer} applied to both the coarse and the fine predictions,
\begin{equation}
  L \;=\; \mathrm{CD}(C, G) \;+\; \mathrm{CD}(F, G),
  \label{eq:routeB_loss}
\end{equation}
where $C$ and $F$ are the coarse and fine outputs. As explained in
\secref{sec:lit_completion}, the symmetric form combines an accuracy term
(predictions must be near the reference) with a completeness term (every
reference point must be covered), preventing collapse to a small safe subset.

\subsection{Optimiser and schedule}
Training uses Adam \citep{kingma2015adam} with learning rate $10^{-3}$ and
batch size 32--48; a StepLR scheduler halves the learning rate every 40 epochs.
All data are preloaded onto the GPU once to avoid per-batch host transfers. On
an NVIDIA RTX~4060 Laptop GPU (8~GB), ten epochs of 6000 pairs take roughly
26--28 minutes. Checkpoints are saved per epoch (for contact sheets) and as a
final \texttt{.pth} file.

\subsection{Training runs}
\label{sec:routeB_runs}
Four training runs were executed:

\begin{itemize}
  \item \texttt{pcn\_dummy\_suop\_10ep.pth} --- dummy, mesh-conditioned
  (Generator~2), 10 epochs;
  \item \texttt{pcn\_drum\_suop\_10ep.pth} --- drum, Generator~1, 10 epochs;
  \item \texttt{pcn\_chair.pth} --- chair, Generator~1, 150 epochs
  (the long run that produced the epoch-1-to-15 progression of
  \figref{fig:chair_pcn_progression});
  \item \texttt{pcn\_dummy\_10ep.pth} / \texttt{pcn\_dummy\_5ep.pth} ---
  exploratory dummy runs used to tune the simulation and grid size.
\end{itemize}

The training progress of the dummy and drum runs is shown in
\figref{fig:dummy_epochs} and \figref{fig:drum_epochs}; the corresponding loss
curves are reported in \secref{ch:resultsB}.

\section{Summary}
\label{sec:routeB_summary}

\routeB{} is chosen when only one view exists and a complete representation is
required. It simulates sonar-like partials from complete references, encodes a
partial with a PointNet-style architecture, and decodes a coarse-to-fine
completion with a folding network, trained end-to-end with a symmetric Chamfer
loss. Its training targets come from procedural models and --- uniquely ---
from the project's own \routeA{} meshes, creating a principled bridge between
the two routes. The next chapter compares the two routes directly and makes
the difference in their outputs explicit.
